% Figure: strain energy density as the triangular area under the linear
% stress-strain line. The left panel shows a small material element under a
% uniaxial stress; the right panel shows the stress-strain line whose
% enclosed triangle is the energy stored per unit volume.
\begin{figure}[htbp]
\centering
\begin{tikzpicture}[
  axis/.style={draw=engblack, -{Stealth}, thick},
  line/.style={draw=engblue, very thick},
  fillU/.style={fill=engblue!16},
  arrow/.style={draw=engred, -{Stealth}, very thick},
  label/.style={font=\small},
]
% left panel: a small cubic element pulled in one direction
\draw[draw=engblack, thick, fill=enggray!10] (0,0.4) rectangle (1.4,1.8);
\draw[arrow] (-0.7,1.1) -- (-0.05,1.1);
\draw[arrow] (2.1,1.1) -- (1.45,1.1);
\node[engred, font=\scriptsize, anchor=south] at (-0.4,1.15) {$\sigma$};
\node[engred, font=\scriptsize, anchor=south] at (1.85,1.15) {$\sigma$};
\node[engblack, font=\scriptsize, anchor=north] at (0.7,0.35) {element $dV$};

% right panel: stress-strain line with triangular area
\begin{scope}[xshift=4.0cm]
\draw[axis] (0,0) -- (4.6,0) node[anchor=west, label] {$\varepsilon$};
\draw[axis] (0,0) -- (0,3.4) node[anchor=south, label] {$\sigma$};
% the linear line sigma = E epsilon up to (eps*, sig*)
\fill[fillU] (0,0) -- (3.6,2.7) -- (3.6,0) -- cycle;
\draw[line] (0,0) -- (3.8,2.85);
\node[engblue, label, anchor=south east] at (3.4,2.3) {$\sigma=E\varepsilon$};
% guide lines to the chosen state
\draw[draw=enggray, dashed, thin] (3.6,0) -- (3.6,2.7);
\draw[draw=enggray, dashed, thin] (0,2.7) -- (3.6,2.7);
\fill[engblack] (3.6,2.7) circle (1.6pt);
\node[font=\scriptsize, anchor=north] at (3.6,-0.05) {$\varepsilon^{*}$};
\node[font=\scriptsize, anchor=east] at (-0.05,2.7) {$\sigma^{*}$};
\node[engblue, font=\scriptsize, anchor=north west] at (1.7,0.55)
  {$u=\tfrac12\sigma^{*}\varepsilon^{*}$};
\end{scope}

\end{tikzpicture}
\caption[Strain energy density as the area under the stress-strain line]{A
small material element (left) carrying a uniaxial stress $\sigma$ stores
elastic energy as it strains. For a linear material the stress-strain line is
straight, $\sigma=E\varepsilon$ (right), and the energy stored per unit volume
is the triangular area beneath it, $u=\tfrac12\sigma^{*}\varepsilon^{*}=
\sigma^{*2}/(2E)$, evaluated at the working state $(\varepsilon^{*},
\sigma^{*})$. Integrating this density over the volume of a loaded member,
with the stress written in terms of the internal load, recovers the
member's total strain energy: $N^2L/(2EA)$ for an axial bar,
$\int M^2/(2EI)\,dx$ for a beam, and $T^2L/(2GJ)$ for a shaft. The
triangular area is the structural sibling of the spring well $\tfrac12kx^2$.}
\label{fig:vol03ch05:strain-energy-density}
\end{figure}
